\section{Finite smooth models and controlled inverses}
\label{sec:inverse}
This section covers every fixed integer $p\ge1$, using the counting theorem
appropriate to that $p$. Its subject is inversion of explicit finite smooth
models. We do not differentiate the discrete sequence, nor an unspecified
error in its asymptotic expansion.

\subsection{Unified notation and model derivatives}
Put
\begin{equation}\label{eq:inverseparams}
 \beta=1/q,\qquad \kappa=p/q,\qquad \alpha=c_p,\qquad
 d_p=\begin{cases}3e/8,&p=1,\\0,&p\ge2,\end{cases}
\end{equation}
and use the smooth positive function
\begin{equation}\label{eq:smoothM}
 M_p(x)=q^{-1}(x/q)^{\kappa x}\exp(\alpha x^\beta-d_p),
 \qquad x>0.
\end{equation}
For a fixed $K\ge0$, define the finite sum
\begin{equation}\label{eq:QK}
 Q_K(x)=\sum_{j=0}^K C_{p,j}(\alpha)x^{-j/q},\qquad
 \nu_K=(K+1)/q.
\end{equation}
For $p=1$, these are the separate critical coefficients from
Theorem~\ref{thm:critical}. Both counting theorems give
\begin{equation}\label{eq:modelvalue}
 a_p(n)/M_p(n)=Q_K(n)+O_{p,K}(n^{-\nu_K}).
\end{equation}
Let
\begin{equation}\label{eq:Fmodel}
 F_K(x)=\kappa x\log(x/q)+\alpha x^\beta-\log q-d_p+\log Q_K(x)
\end{equation}
where $x$ is sufficiently large that $Q_K(x)>0$.

\begin{lemma}\label{lem:model}
For each fixed $p,K$, there exists $X_{p,K}>e$ such that for
$x\ge X_{p,K}$,
\begin{equation}\label{eq:modelderivative}
 Q_K(x)>1/2,\qquad F_K'(x)\ge c_p'\log x>0,\qquad
             0<F_K''(x)\le C_{p,K}/x.
\end{equation}
The function $F_K$ has a unique smooth inverse $u_K$ on its eventual range,
and at integers
\begin{equation}\label{eq:logcomparison}
 \log a_p(n)=F_K(n)+O_{p,K}(n^{-\nu_K}).
\end{equation}
For fixed $L\ge K$ and each fixed derivative order $h\ge0$,
\begin{equation}\label{eq:coherentF}
 (F_L-F_K)^{(h)}(x)=O_{p,K,L,h}(x^{-\nu_K-h}).
\end{equation}
\end{lemma}
\begin{proof}
Let $\delta=1/2$ for $p=1$ and $\delta=(p-1)/q$ for $p\ge2$.
The vanishing coefficients already proved imply
$(Q_K-1)^{(h)}(x)=O_{p,K,h}(x^{-\delta-h})$;
if $Q_K=1$ the left side is zero. This is ordinary differentiation of a
known finite expression. In particular $Q_K>1/2$ eventually, and
$(\log Q_K)^{(h)}=O(x^{-\delta-h})$.
Thus
\begin{align}
 F_K'(x)&=\kappa(\log(x/q)+1)+\alpha\beta x^{\beta-1}
                                      +O(x^{-\delta-1}),\label{eq:Fprime}\\
 F_K''(x)&=\kappa/x+\alpha\beta(\beta-1)x^{\beta-2}
                                      +O(x^{-\delta-2}).\label{eq:Fsecondmodel}
\end{align}
The positive $\kappa/x$ term dominates the second derivative, proving
\eqref{eq:modelderivative}. Higher derivatives of order $h\ge2$ are
$O_{p,K,h}(x^{1-h})$. Taking a logarithm in \eqref{eq:modelvalue}, using
$Q_K>1/2$, proves \eqref{eq:logcomparison}.
Finally $Q_L-Q_K$ is a finite sum starting at $x^{-\nu_K}$, so its
fixed derivatives and those of the logarithmic difference have the stated
bounds. This proves \eqref{eq:coherentF}.
\end{proof}

\subsection{A Lambert core and finitely many Newton steps}
Write $Y=\log y$ and set
\begin{equation}\label{eq:Lambertcore}
 Z=Y+\log q+d_p,\qquad
 N=N(Y)=\frac{qZ}{pW(Z/p)},
\end{equation}
where $W$ is the positive real branch.
This solves exactly
\begin{equation}\label{eq:coreeq}
 \kappa N\log(N/q)=Z.
\end{equation}
Indeed, $N/q=\exp(W(Z/p))$ and $W(Z/p)e^{W(Z/p)}=Z/p$.
For fixed $K,J$, begin at $x_0=N(Y)$ and make $J$ Newton steps
\begin{equation}\label{eq:newton}
 x_{j+1}=x_j-\frac{F_K(x_j)-Y}{F_K'(x_j)},\qquad 0\le j<J.
\end{equation}
This is an explicit finite composition of $W$, real powers, logarithms,
exponentials, and arithmetic. There is no hidden root solve in $x_J$.

\begin{theorem}\label{thm:newton}
For every fixed $p,K,J$, all iterates in \eqref{eq:newton} are well-defined
for sufficiently large $Y$, and
\begin{equation}\label{eq:monotoneNewton}
 u_K(Y)\le x_{j+1}\le x_j\le N(Y).
\end{equation}
Their error obeys
\begin{equation}\label{eq:newtonerror}
 0\le x_J-u_K(Y)
   \le C_{p,K,J}\,
       \frac{N^{1-(1-\beta)2^J}}{(\log N)^{2^{J+1}-1}}.
\end{equation}
\end{theorem}
\begin{proof}
At the starting point,
$F_K(N)-Y=\alpha N^\beta+\log Q_K(N)>0$ eventually.
On the other hand $F_K(N/2)<Y$ for large $Y$, because the change in the
core term is of order $N\log N$, much larger than $N^\beta$.
Thus $u=u_K(Y)\in(N/2,N)$.
By the mean-value theorem and \eqref{eq:modelderivative},
\begin{equation}\label{eq:e0}
 e_0=N-u=O_{p,K}(N^\beta/\log N).
\end{equation}
On $[u,N]$, $F_K$ is increasing and convex. The tangent line at $x\ge u$
meets the line of height $Y$ between $u$ and $x$; this proves the
monotonicity and all domain assertions in \eqref{eq:monotoneNewton}.
Taylor's theorem, expanded at $x_j$, gives
\[
 e_{j+1}=\frac{F_K''(\xi_j)}{2F_K'(x_j)}e_j^2
                \le C\frac{e_j^2}{N\log N}.
\]
The power of $N$ therefore changes by $a\mapsto2a-1$ and the power in
the denominator by $b\mapsto2b+1$. Starting from $(a,b)=(\beta,1)$ in
\eqref{eq:e0} yields \eqref{eq:newtonerror}.
\end{proof}

The first Newton error is $O(N^{2\beta-1}/\log^3N)$.
At criticality $\beta=1/2$, so the first two Newton errors are respectively
\begin{equation}\label{eq:criticalNewton}
 O(\log^{-3}N),\qquad O(N^{-1}\log^{-7}N).
\end{equation}
The constant offset in the critical core is essential:
$N=2(Y+\log2+3e/8)/W(Y+\log2+3e/8)$.

\subsection{Arbitrary fixed index accuracy}
\begin{corollary}\label{cor:index}
Given any fixed $M>0$, choose integers $K,J\ge0$ with
\begin{equation}\label{eq:chooseKJ}
 (K+1)/q\ge M,\qquad (1-\beta)2^J\ge M+1.
\end{equation}
Let $h_{K,J}(Y)=x_J(Y)$. Then on the exact range of the count sequence,
\begin{equation}\label{eq:indexrecover}
 h_{K,J}(\log a_p(n))=n+O_{p,K,J}(n^{-M}/\log n).
\end{equation}
In particular, its nearest integer is $n$ for all sufficiently large $n$.
\end{corollary}
\begin{proof}
By \eqref{eq:logcomparison} and the derivative lower bound,
\[
 u_K(\log a_p(n))-n=O_{p,K}(n^{-\nu_K}/\log n).
\]
The leading growth also gives $N(\log a_p(n))\sim n$.
Combine this comparison with \eqref{eq:newtonerror} and
\eqref{eq:chooseKJ}. The logarithmic denominator in the Newton bound
is at least $\log N$, giving \eqref{eq:indexrecover}.
Eventually its absolute error is less than $1/2$.
\end{proof}

The inverse models are coherent. From \eqref{eq:coherentF} and the
mean-value theorem, for each fixed $L\ge K$,
\begin{equation}\label{eq:inversecoherent}
 u_L(Y)-u_K(Y)=O_{p,K,L}(N(Y)^{-\nu_K}/\log N(Y)).
\end{equation}
Increasing the model order therefore cannot change already resolved inverse
terms. All derivative bounds used here belong to explicit finite functions;
none is inferred by differentiating \eqref{eq:logcomparison}.

\subsection{A convergent local Lagrange generator}
For a fixed sufficiently large $N=N(Y)$, put
\begin{equation}\label{eq:HN}
 v=F_K(N)-Y,\qquad
 H_N(z)=\frac{z}{F_K(N+z)-F_K(N)},\qquad H_N(0)=\frac1{F_K'(N)}.
\end{equation}
The local Lagrange formula is
\begin{equation}\label{eq:Lagrange}
 u_K(Y)=N+\sum_{m\ge1}\frac{(-v)^m}{m!}
       \left.\frac{d^{m-1}}{dz^{m-1}}H_N(z)^m\right|_{z=0}.
\end{equation}
We verify a convergence disk and a quantitative truncation bound, rather
than using this as a merely formal identity.

\begin{proposition}\label{prop:Lagrange}
For fixed $p,K$ and all sufficiently large $Y$, series
\eqref{eq:Lagrange} converges. After retaining $m=1,\ldots,L$ its remainder is
\begin{equation}\label{eq:Lagrangerem}
 O_{p,K,L}\left(N\left(\frac{N^{\beta-1}}{\log N}\right)^{L+1}\right).
\end{equation}
\end{proposition}
\begin{proof}
Fix $0<\eta<1/4$. The explicit finite expression $F_K(N+z)$ is holomorphic
on a neighborhood of $|z|\le\eta N$ for large $N$: $N+z$ stays in the
right half-plane, and $Q_K(N+z)$ stays near one.
Uniformly on this disk,
\[
 F_K'(N+z)=\kappa\log N+O_{p,K,\eta}(1).
\]
Its real part is positive eventually. For two points in the disk, divide
their function-value difference by their point difference and integrate
the derivative along their connecting segment. The real part of the
result is positive, so the map is injective.

Write $G(z)=F_K(N+z)-F_K(N)$. Integrating the same derivative bound gives
\[
 |G(z)-\kappa(\log N)z|\le C|z|.
\]
Let $R_N=(\kappa\eta/4)N\log N$.
For large $N$, on $|z|=\eta N$ the perturbation is at most
$\tfrac14\kappa\eta N\log N$. For every $|w|<2R_N$,
\[
 |G(z)-\kappa(\log N)z-w|
      <\tfrac34\kappa\eta N\log N
      <|\kappa(\log N)z|.
\]
Rouch\'e's theorem gives exactly one solution of $G(z)=w$ inside this disk.
The derivative is nonzero there, so the inverse displacement is analytic
on the open image disk $|w|<2R_N$ and bounded by $\eta N$.
This strict radius margin permits Cauchy's estimate on $|w|=R_N$.

The Taylor coefficients of the inverse displacement are, by the local
Lagrange theorem~\cite{dlmf},
$\frac1{m!}(d^{m-1}/dz^{m-1})H_N(z)^m|_{z=0}$.
Since $v=O(N^\beta)$, one has $|v|/R_N\to0$.
Cauchy's bound and a geometric tail yield
\[
 O\left(N\frac{(|v|/R_N)^{L+1}}{1-|v|/R_N}\right),
\]
which proves \eqref{eq:Lagrangerem} and convergence at $w=-v$.
\end{proof}

The first two displacements in \eqref{eq:Lagrange} are
\begin{equation}\label{eq:firstLagrange}
 -\frac v{F_K'(N)}-\frac{F_K''(N)v^2}{2F_K'(N)^3}.
\end{equation}
The first leading displacement alone is
\begin{equation}\label{eq:firstdisplacement}
 -\frac{\alpha N^\beta}{\kappa(\log(N/q)+1)}.
\end{equation}
Discarding terms to get \eqref{eq:firstdisplacement} also discards their
accuracy. This shorter expression is not assigned the much smaller error
of a multistep Newton approximation.

\subsection{Exact integer thresholds and the rounding boundary}
\begin{theorem}\label{thm:threshold}
Fix $M>0$ and $K,J$ satisfying \eqref{eq:chooseKJ}.
For $h=h_{K,J}(\log y)$ there are constants $C,Y_0$, depending only on
$p,K,J$, such that for $y\ge e^{Y_0}$,
\begin{equation}\label{eq:ceiling}
 \left\lceil h-C\frac{h^{-M}}{\log h}\right\rceil
       \le T_p(y)\le
 \left\lceil h+C\frac{h^{-M}}{\log h}\right\rceil.
\end{equation}
The endpoint integers differ by at most one for all sufficiently large $y$.
If they coincide, that integer is the exact threshold. Otherwise the theorem
leaves two adjacent possibilities.
\end{theorem}
\begin{proof}
First use the exact model inverse $u=u_K(Y)$, $Y=\log y$.
For integers $m$ with $|m-u|\le2$, the error in
\eqref{eq:logcomparison} is bounded by $C_1u^{-\nu_K}$ and
$F_K'(x)\ge c_1\log u$ throughout this small neighborhood.
Choose
\[
 \varepsilon=\frac{2C_1u^{-\nu_K}}{c_1\log u}.
\]
For $m_+=\lceil u+\varepsilon\rceil$,
$F_K(m_+)-Y\ge2C_1u^{-\nu_K}$, hence
$\log a_p(m_+)\ge Y$.
For $m_-=\lceil u-\varepsilon\rceil-1$ one has strictly
$m_-<u-\varepsilon$, even if $u-\varepsilon$ is itself an integer.
It follows that $\log a_p(m_-)<Y$.
Exact monotonicity from Proposition~\ref{prop:monotone} yields
\[
 \lceil u-\varepsilon\rceil\le T_p(y)
                         \le\lceil u+\varepsilon\rceil.
\]
Replacing $u$ by $h$ uses \eqref{eq:newtonerror}; $u\sim h\sim N$ and
$\nu_K\ge M$ absorb both errors into the radius in \eqref{eq:ceiling}.
Twice that radius is eventually less than one, so its endpoint ceilings
can differ by at most one.
\end{proof}

\begin{remark}[Exact range versus arbitrary levels]
Nearest-integer recovery in Corollary~\ref{cor:index} assumes the input
is exactly $y=a_p(n)$. It does not give nearest-integer threshold recovery
for an arbitrary $y$. Even an extremely short real interval can straddle
an integer, in which case taking a ceiling cannot be settled by its width
alone. If actual numerical bounds and candidates are available, direct
integer evaluation can decide the remaining threshold comparison.
The present asymptotic theorem does not supply certified finite constants
or a certified onset for doing so.
\end{remark}

All onset thresholds and error constants in this inverse section are
existential. They could in principle be made effective from explicit
Taylor and concentration bounds, but the accompanying numerical inverse
experiments are diagnostics, not a certified finite-input algorithm.
The critical count theorem transfers to these inverses because its finite
models have the same derivative structure. It does not transfer the marked
probability theorem to $p=1$.
